\chapter{Accessibility Metric System}

To provide a basis of the subjective development experience, each source will be subjected to an analysis which tracks the frequency
and severity of instances that the source is lacking in clarity and requires switching to other resources to understand, or add to,
the context written.

\section{Scoring Stystem}

Each source will accumulate points depending on how far and how long I need to refer to other sources to understand the original.

\begin{table}[h]
\centering
\caption{Complexity Severity Levels}
\label{tab:severity}
\begin{tabularx}{\textwidth}{@{} ll X l @{}}
\toprule
\textbf{Level} & \textbf{Value} & \textbf{Description} & \textbf{Examples} \\ \midrule
Minor    & 1 & Lexical issue. A word, phrase, or term that may not be immediately understood, requiring brief additional research. & 'Isospin', 'Alpha Decay' \\ \addlinespace
Moderate & 2 & Notation barrier. Mathematical notation that is not explained, requiring research into determining its meaning. & $\sum$ \\ \addlinespace
Major    & 4 & Unexplained introduced concepts. Concepts that are mentioned but not elaborated, which may require extensive further research. & --- \\ \addlinespace
Critical & 8 & Incomprehensive topic. A topic is difficult to research or understand in the time given; e.g., a proof with no numerical constants. & Complex proofs \\ \bottomrule
\end{tabularx}
\end{table}

\subsection{Evaluation Process}

As the particle simulation is developed, each time the workflow is interrupted with needing to find context through
additional resources, points will be accumulated and logged for each major resource, and a final 
\textbf{Approachability Score} will be created using the following equation:

\begin{equation}
    A = \sum (P_{\text{minor}} + P_{\text{moderate}} + P_{\text{major}} + P_{\text{critical}})
\end{equation}

A lower Approachability Score indicates a more developer-friendly resource. Whereas a higher score suggests a more difficult resource
for developers to use.